\documentclass[12pt,a4paper]{article}

\usepackage[utf8]{inputenc}
\usepackage[T1]{fontenc}
\usepackage[spanish]{babel}
\usepackage{lmodern}
\usepackage{geometry}

\geometry{top=2.5cm,bottom=2.5cm,left=2.5cm,right=2.5cm}

\setlength{\parindent}{0pt}
\setlength{\parskip}{0.6em}

\begin{document}

\begin{center}
	{\Large\bfseries MEMORIA JUSTIFICATIVA}
\end{center}

Yo, Pablo de Ferra Chermaz, estudiante del Master en Astrofisica, expongo mediante la presente la justificacion de mi desplazamiento al Barcelona Supercomputing Center (BSC-CNS), realizado durante los dias4 a 6 de diciembre de 2025.

El objetivo principal de la visita fue conocer de primera mano las infraestructuras y capacidades de supercomputacion del BSC, asi como su aplicacion a problemas de investigacion y computacion cientifica. Esta actividad contribuye directamente a mi formacion academica y a la adquisicion de competencias relacionadas con el uso de recursos HPC (High Performance Computing) en el contexto del Master.

Durante la estancia se realizaron actividades academicas de caracter formativo, incluyendo visitas a las instalaciones y sesiones de divulgacion tecnica sobre los recursos, servicios y lineas de investigacion del centro (p. ej., arquitectura de sistemas, planificacion de trabajos, almacenamiento y buenas practicas de uso de supercomputacion).

La visita me permitio comprender el entorno de trabajo asociado a grandes infraestructuras de calculo, familiarizarme con la terminologia y flujos habituales (colas, asignacion de recursos, ejecucion reproducible, etc.) y contextualizar el uso de HPC en problemas actuales de astrofisica y computacion cientifica. Asimismo, la actividad refuerza mi preparacion para el desarrollo de trabajos que requieran recursos computacionales avanzados.

En La Laguna, a 20 de Marzo de 2026.

\vspace{2.2cm}

\textbf{Fdo.:} Pablo de Ferra Chermaz

\end{document}
